\section{Objectifs}

\begin{itemize}
\item Déterminer une équation de droite à partir de données.
\item Lire une pente ou un vecteur directeur.
\item Tracer une droite à partir de son équation.
\item Déterminer le point d'intersection de deux droites sécantes.
\end{itemize}

\section{Prérequis}

\begin{itemize}
\item fonctions affines
\item vecteurs
\item résolution d'équations
\end{itemize}

\section{Activité de découverte}

Deux forfaits de téléphone sont représentés par deux droites. Leur intersection indique le volume de données pour lequel les prix sont égaux.

\section{Vocabulaire}

\begin{itemize}
\item \textbf{droite}
\item \textbf{équation réduite}
\item \textbf{équation cartésienne}
\item \textbf{pente}
\item \textbf{coefficient directeur}
\item \textbf{vecteur directeur}
\item \textbf{sécantes}
\item \textbf{parallèles}
\end{itemize}

\section{Cours}

\subsection{Définitions}

\begin{itemize}
\item Une équation réduite de droite non verticale s'écrit y = mx + p.
\item m est la pente ou coefficient directeur et p l'ordonnée à l'origine.
\item Une équation cartésienne de droite s'écrit ax + by + c = 0 avec a et b non tous deux nuls.
\end{itemize}

\subsection{Propriétés}

\begin{itemize}
\item Deux droites non verticales de même pente sont parallèles.
\item Un vecteur directeur d'une droite ax + by + c = 0 est (-b ; a).
\item Le point d'intersection de deux droites sécantes vérifie les deux équations.
\end{itemize}

\subsection{Méthodes}

\begin{enumerate}
\item Avec deux points A et B, calculer la pente m=(yB-yA)/(xB-xA) si xA≠xB.
\item Pour tracer y=mx+p, placer l'ordonnée à l'origine puis utiliser la pente.
\item Pour une intersection, résoudre le système formé par les deux équations.
\end{enumerate}

\section{Exemples corrigés}

\begin{itemize}
\item La droite passant par A(0 ; 1) et B(2 ; 5) a pour pente 2, donc y=2x+1.
\item Les droites y=3x-1 et y=3x+4 sont parallèles car elles ont la même pente.
\end{itemize}

\coursefigure{/images/mathematiques/latex-migration/2nde-droites-plan-figure-1.svg}{Deux droites se coupent en un point d'intersection.}{Droites du plan : représentation schématique à commenter avant les exercices.}

\section{Algorithmique en Python}

\subsection{Équation réduite avec deux points}

\begin{verbatim}
def equation_reduite(A, B):
    m = (B[1] - A[1]) / (B[0] - A[0])
    p = A[1] - m * A[0]
    return m, p

print(equation_reduite((0, 1), (2, 5)))
\end{verbatim}

On calcule d'abord la pente, puis p à partir d'un point de la droite.

\section{Erreurs fréquentes}

\begin{itemize}
\item Utiliser la formule de la pente pour une droite verticale.
\item Confondre pente et ordonnée à l'origine.
\item Dire que deux droites de pentes différentes sont parallèles.
\end{itemize}

\section{Formules et idées à retenir}

\begin{itemize}
\item y = mx + p
\item m = (yB-yA)/(xB-xA)
\item ax + by + c = 0
\item vecteur directeur (-b ; a)
\end{itemize}

\section{Synthèse}

Ce chapitre sert à installer des automatismes, mais aussi à choisir une représentation adaptée : texte, formule, tableau, figure, graphique ou programme. Les exercices associés reprennent les cinq niveaux attendus : automatismes, application directe, méthode guidée, raisonnement et synthèse.
